% The recording deck as a Beamer presentation, built on the same skeleton
% as the Susquehanna Algothon deck (Madrid theme, one palette pass, the
% sig* macro layer) so slides port between the two. Compile with
% `make slides` (tectonic presentation/deck.tex). Asset paths are relative
% to the presentation/ directory, so the file has to stay where it lives.
%
% The palette values come from scripts/figstyle.py, so the slides and the
% figures read as one set. The MP4 animations cannot play inside a PDF;
% those slides carry a still frame from presentation/stills/ (make stills)
% and the moving version is spliced in fullscreen during the edit, as
% presentation/SHOTLIST.md prescribes. Code appears as screenshots of the
% real source (make code-shots), never retyped listings, so every frame
% carries its file:line provenance.

\documentclass[aspectratio=169,11pt]{beamer}

\usepackage[T1]{fontenc}
\usepackage{lmodern}
\usepackage{microtype}
\usepackage{xstring}
\usepackage{amsmath}
\usepackage{bm}
\usepackage{graphicx}
\usepackage{booktabs}
\usepackage{colortbl}
\usepackage{tikz}

% ***************************************************
% Palette
% ***************************************************
% figstyle.py values wearing the sig* role names, so slide code written
% for either deck compiles against either palette.

\definecolor{sigblue}{HTML}{2563EB}   % BLUE, the thing that worked
\definecolor{sigink}{HTML}{1F2937}    % deck.html --ink, titles and text
\definecolor{signavy}{HTML}{111827}   % darker ink for the head and foot bars
\definecolor{sigslate}{HTML}{6B7280}  % DARK_GREY, captions and notes
\definecolor{sigmist}{HTML}{F0F3F9}   % background canvas, a shade of wash
\definecolor{sigcard}{HTML}{FFFFFF}   % figure and block bodies
\definecolor{sigline}{HTML}{D1D5DB}   % SPINE, hairlines
\definecolor{sigwhite}{HTML}{FFFFFF}
\definecolor{siggreen}{HTML}{059669}  % example blocks
\definecolor{sigamber}{HTML}{D97706}  % AMBER, alert blocks
\definecolor{sigred}{HTML}{DC2626}    % RED, negative controls

% ***************************************************
% Theme
% ***************************************************

\usetheme{Madrid}
\usecolortheme{default}
\usefonttheme{professionalfonts}
\usefonttheme[onlymath]{serif}

\setbeamercolor{background canvas}{bg=sigmist}

\setbeamercolor*{palette primary}{bg=signavy, fg=sigwhite}
\setbeamercolor*{palette secondary}{bg=sigblue, fg=sigwhite}
\setbeamercolor*{palette tertiary}{bg=sigink, fg=sigwhite}
\setbeamercolor*{palette quaternary}{bg=signavy, fg=sigwhite}

% Every beamer element that carries a palette colour, assigned in one pass
% so the theme has a single source of truth.
\makeatletter
\@for\@entry:={%
  normal text/fg/sigink,structure/fg/sigblue,%
  itemize item/fg/sigblue,itemize subitem/fg/sigslate,%
  enumerate item/fg/sigblue,alerted text/fg/sigblue,caption name/fg/sigblue,%
  section in toc/fg/sigink,subsection in toc/fg/sigslate,%
  block title/bg/sigblue,block title/fg/sigwhite,%
  block title alerted/bg/sigamber,block title alerted/fg/sigwhite,%
  block title example/bg/siggreen,block title example/fg/sigwhite,%
  block body/bg/sigcard,block body/fg/sigink,%
  block body alerted/bg/sigcard,block body alerted/fg/sigink,%
  block body example/bg/sigcard,block body example/fg/sigink%
}\do{%
  \StrCut{\@entry}{/}\@element\@rest
  \StrCut{\@rest}{/}\@channel\@value
  % Beamer stores the key list unexpanded, so the assignment is built now.
  \edef\@assign{\noexpand\setbeamercolor{\@element}{\@channel=\@value}}%
  \@assign
}
\makeatother

\setbeamerfont{author in head/foot}{size=\tiny}
\setbeamertemplate{footline}{%
  \leavevmode%
  \hbox{%
    \begin{beamercolorbox}[wd=.44\paperwidth,ht=2.25ex,dp=1ex,center]{author in head/foot}%
      \usebeamerfont{author in head/foot}\insertshortauthor
    \end{beamercolorbox}%
    \begin{beamercolorbox}[wd=.31\paperwidth,ht=2.25ex,dp=1ex,center]{title in head/foot}%
      \usebeamerfont{title in head/foot}\insertshorttitle
    \end{beamercolorbox}%
    \begin{beamercolorbox}[wd=.25\paperwidth,ht=2.25ex,dp=1ex,right]{date in head/foot}%
      \usebeamerfont{date in head/foot}\insertshortdate{}\hspace*{2em}%
      \insertframenumber{} / \inserttotalframenumber\hspace*{2ex}%
    \end{beamercolorbox}}%
  \vskip0pt%
}

\setbeamertemplate{navigation symbols}{}
\setbeamertemplate{itemize item}{\raisebox{0.15ex}{\rule{0.6ex}{0.6ex}}}
\setbeamertemplate{itemize subitem}{\raisebox{0.2ex}{\rule{0.45ex}{0.45ex}}}
\setbeamertemplate{blocks}[rounded][shadow=false]

% ***************************************************
% Furniture
% ***************************************************

\newcommand{\sigemph}[1]{{\color{sigblue}\bfseries #1}}

% ***************************************************
% Slide structures
% ***************************************************

% Row of figures with captions beneath. Column width follows the number of
% entries. Values must not contain a slash, it is the pair separator.
\newcounter{sigentries}
\newcommand{\sigstats}[1]{%
  \setcounter{sigentries}{0}%
  \foreach \sigignored in {#1}{\stepcounter{sigentries}}%
  \pgfmathsetmacro{\sigcolwidth}{(1 - 0.04*(\value{sigentries}-1))/\value{sigentries}}%
  \par\medskip\noindent
  \foreach \sigvalue/\sigcaption [count=\sigindex] in {#1}{%
    \ifnum\sigindex>1\hfill\fi
    \begin{minipage}[t]{\sigcolwidth\linewidth}
      \raggedright{\color{sigblue}\Large\bfseries \sigvalue}\par\vspace{0.15em}
      {\color{sigslate}\scriptsize \sigcaption}
    \end{minipage}%
  }%
  \par
}

% Tinted band across the foot of a slide, for one sentence.
\newcommand{\sigcallout}[1]{%
  \vfill
  \begin{beamercolorbox}[wd=\linewidth,sep=0.7em,rounded=true]{block body}
    {\color{sigblue}\rule{0.9ex}{0.9ex}}\hspace{0.5em}#1
  \end{beamercolorbox}%
}

% Two plain dot points under a rung's screenshot. The first says what the
% technique changes, the second why it was expected to pay. The measured
% verdict stays in the callout beneath them.
\newcommand{\sigrungnotes}[2]{%
  \par\vspace{0.25em}
  {\footnotesize
  \begin{itemize}\setlength{\itemsep}{0.1em}
    \item #1
    \item #2
  \end{itemize}}%
}

% Balanced two-column body, gutter set once.
\newenvironment{sigcolumns}{\begin{columns}[T,onlytextwidth]}{\end{columns}}
\newenvironment{sigcolumn}{\begin{column}{0.48\textwidth}}{\end{column}}

% Booktabs table with the rules, sizing and stretch set once.
\newcommand{\sigtable}[3]{%
  \begingroup
  \renewcommand{\arraystretch}{1.15}%
  \begin{table}\small
    \begin{tabular}{#1}
      \toprule
      #2 \\
      \midrule
      #3
      \bottomrule
    \end{tabular}
  \end{table}%
  \endgroup
}

% Figure slot. Takes the rendered asset's path; until the asset exists the
% slot reserves the same space and names the file it wants, so the layout
% can be judged before make stills / make code-shots have run.
\newlength{\sigslotheight}
\setlength{\sigslotheight}{0.54\textheight}
% Setting \sigslotheight inside a frame stays local to that frame, which is
% how the rung frames trade figure height for their dot points.
\newcommand{\sigrungfigure}{\setlength{\sigslotheight}{0.34\textheight}}
\newcommand{\sigfigure}[3][\linewidth]{%
  \begin{minipage}{#1}
    \centering
    \IfFileExists{#2}{%
      \includegraphics[width=\linewidth,height=\sigslotheight,keepaspectratio]{#2}%
    }{%
      \setlength{\fboxsep}{0pt}%
      \fcolorbox{sigline}{sigcard}{%
        \begin{minipage}[c][\sigslotheight][c]{\dimexpr\linewidth-2\fboxrule\relax}
          \centering
          {\color{sigslate}\scriptsize\texttt{#2}}
        \end{minipage}%
      }%
    }%
    \par\vspace{0.45em}
    {\scriptsize\color{sigslate} #3\par}
  \end{minipage}%
}

% Supporting frames sit after the appendix marker and outside the count.
\newcounter{siglastframe}
\newcommand{\sigbackupbegin}{\setcounter{siglastframe}{\value{framenumber}}}
\newcommand{\sigbackupend}{\setcounter{framenumber}{\value{siglastframe}}}

% ***************************************************
% Title data
% ***************************************************

\title[Heston PDE pricer]{Pricing an option by solving a PDE}
\subtitle{Serial implementation, optimisation and benchmarking}
\author[48008361]{Student 48008361}
\institute[]{COSC3500 Milestone 1}
\date{August 2026}

\begin{document}

\begin{frame}[plain]
  \titlepage
\end{frame}

% --------------------------------------------------------------------------
% Introduction. The assessors are HPC staff, not finance people, so this
% section builds the problem from zero.
% --------------------------------------------------------------------------

\begin{frame}{What an option is}{The right, but not the obligation, to buy at a fixed price}
  \begin{sigcolumns}
    \begin{sigcolumn}
      \begin{itemize}
        \item A European call is the right, but not the obligation, to buy
          a share for a fixed \sigemph{strike} price on a fixed day.
        \item At expiry its value is known exactly. Above the strike it
          pays the difference, below it pays nothing, so the payoff has a
          kink at the strike.
        \item The problem is the price \sigemph{today}, ninety-one days
          before expiry.
      \end{itemize}
    \end{sigcolumn}
    \begin{sigcolumn}
      \sigfigure{../results/payoff.png}{The payoff at expiry, with the
        kink at the strike.}
    \end{sigcolumn}
  \end{sigcolumns}
  \sigstats{5200/today's share price, 5250/the strike, 91 days/to expiry}
\end{frame}

\begin{frame}{Why not Black-Scholes}{Implied volatility is not constant across strikes}
  \begin{sigcolumns}
    \begin{sigcolumn}
      \begin{itemize}
        \item Black-Scholes assumes the share's volatility is
          \sigemph{one constant number}.
        \item Implied volatility computed from market option prices varies
          across strikes, which contradicts that assumption.
        \item The fix is a model in which the volatility itself varies.
      \end{itemize}
    \end{sigcolumn}
    \begin{sigcolumn}
      \sigfigure{stills/still_smile.png}{Implied volatility from Heston
        PDE prices. The dashed line is the single Black-Scholes
        volatility.}
    \end{sigcolumn}
  \end{sigcolumns}
  \sigcallout{Setting Heston's volatility of volatility to zero recovers
    Black-Scholes exactly. That is the main validation test.}
\end{frame}

\begin{frame}{The Heston model}{Two correlated stochastic processes}
  \begin{block}{The share price and its variance are both stochastic processes}
    \vspace{-0.4em}
    \begin{align*}
      dS &= (r - q)\,S\,dt + \sqrt{v}\,S\,dW_1\\[0.1em]
      dv &= \kappa(\theta - v)\,dt + \xi\sqrt{v}\,dW_2,
        \qquad \mathrm{corr}(dW_1, dW_2) = \rho
    \end{align*}
    \vspace{-0.9em}
  \end{block}
  \begin{itemize}
    \item $S$ is the share price and $v$ is its instantaneous variance,
      which follows its own stochastic process.
    \item $\kappa$ is the rate at which the variance reverts to its
      long-run level $\theta$, and $\xi$ sets the size of the variance
      fluctuations.
    \item The two noises are correlated at $\rho = -0.70$.
  \end{itemize}
  \sigcallout{The correlation is negative because volatility tends to rise
    when prices fall.}
\end{frame}

\begin{frame}{The symbols in the two equations}{Definitions}
  \sigtable{cl}{\textbf{Symbol} & \textbf{Definition}}{%
    $S$ & the share price \\
    $v$ & the variance of the share price, the square of its volatility \\
    $r$ & the risk-free interest rate \\
    $q$ & the dividend yield of the share \\
    $dt$ & a small step forward in time \\
    $dW_1,\ dW_2$ & random increments, one for the price and one for the variance \\
    $\kappa$ & the rate at which the variance reverts to its long-run level \\
    $\theta$ & the long-run level of the variance \\
    $\xi$ & the volatility of volatility, which sets the size of variance fluctuations \\
    $\rho$ & the correlation between the two increments \\}
\end{frame}

\begin{frame}{Where the strike price enters}{The model and the contract are separate}
  \begin{itemize}
    \item The two equations describe how the share price moves. They do
      not involve the option contract, which is why no strike appears in
      them.
    \item The contract enters at expiry, where the grid is filled with the
      payoff, $\max(S - K,\, 0)$ for this call. That is where the strike
      $K$ enters.
    \item Black-Scholes merges the model and the contract into one closed
      formula, which is why $S$ and $K$ appear together there. The PDE
      route keeps them separate.
  \end{itemize}
  \sigcallout{Heston does not build on the Black-Scholes formula. It
    relaxes one Black-Scholes assumption, the constant volatility. With
    $\xi = 0$ the model reduces to Black-Scholes, which is the main
    validation test.}
\end{frame}

\begin{frame}{From stochastic model to deterministic equation}{There is no random number generator in the program}
  \begin{itemize}
    \item The fair price is an average. It is the expected payoff over
      every path the market could take, discounted back to today.
    \item One way to compute an average is to simulate many random paths
      and take their mean. This program does not do that.
    \item The Feynman-Kac theorem says that average, viewed as a function
      of the starting point, solves a deterministic equation. That
      function is the option price, written $V$. The following slides
      construct its equation.
  \end{itemize}
  \sigcallout{Averaging over the noise leaves a deterministic equation, so
    the program needs no random numbers.}
\end{frame}

\begin{frame}{The derivative notation}{Subscripts denote partial derivatives}
  \begin{itemize}
    \item $V(S, v, t)$ is the option price when the share price is $S$,
      the variance is $v$ and the time is $t$. It is the unknown the
      program solves for.
    \item Capital $V$ and lowercase $v$ are different quantities. $V$ is
      the option price, $v$ is the variance of the share price.
    \item A subscript denotes a partial derivative, taken with the other
      variables held fixed.
  \end{itemize}
  \sigtable{cl}{\textbf{Symbol} & \textbf{Definition}}{%
    $V_S$ & the first derivative with respect to the share price \\
    $V_{SS}$ & the second derivative with respect to the share price \\
    $V_v$ & the first derivative with respect to the variance \\
    $V_{vv}$ & the second derivative with respect to the variance \\
    $V_{Sv}$ & the mixed derivative, in the share price and the variance \\}
  \sigcallout{On the grid each of these becomes a difference between
    neighbouring cells, which is why one cell update reads nine neighbours.}
\end{frame}

\begin{frame}{Why the noise terms become second derivatives}{Averaging over one time step}
  \begin{itemize}
    \item Over one small time step the noise moves the share price up or
      down by about $\sqrt{v}\,S$, with both directions equally likely.
    \item Averaging the value over the up move and the down move cancels
      the first-derivative contributions.
    \item The remaining term is second order, half the move size squared
      times $V_{SS}$.
  \end{itemize}
  \begin{center}
    move size $\sqrt{v}\,S$
    \qquad so the noise in $S$ becomes \qquad
    $\tfrac{1}{2}\,(\sqrt{v}\,S)^2\, V_{SS} \;=\; \tfrac{1}{2}\,v S^2\, V_{SS}$
  \end{center}
  \sigcallout{The same argument applied to the variance noise, of size
    $\xi\sqrt{v}$, gives $\tfrac{1}{2}\,\xi^2 v\, V_{vv}$. The correlation
    between the two noises gives the mixed term.}
\end{frame}

\begin{frame}{The equation, term by term}{Each part of the model contributes one term}
  \sigtable{ll}{\textbf{Part of the model} & \textbf{Resulting term}}{%
    time passing & $V_t$ \\
    the drift of $S$, which is $(r - q)\,S$ & $(r - q)\,S\,V_S$ \\
    the drift of $v$, which is $\kappa(\theta - v)$ & $\kappa(\theta - v)\,V_v$ \\
    the noise in $S$, of size $\sqrt{v}\,S$ & $\tfrac{1}{2}\,v S^2\, V_{SS}$ \\
    the noise in $v$, of size $\xi\sqrt{v}$ & $\tfrac{1}{2}\,\xi^2 v\, V_{vv}$ \\
    the correlation $\rho$ between the noises & $\rho\,\xi\,v S\, V_{Sv}$ \\
    discounting at the rate $r$ & $r\,V$ \\}
  \sigcallout{Drift terms give first derivatives and noise terms give
    second derivatives. The drift terms need no averaging because they are
    not random.}
\end{frame}

\begin{frame}{The equation the program solves}{Every quantity except $V$ is known}
  \begin{block}{The Heston pricing equation}
    \vspace{-0.4em}
    \begin{equation*}
      V_t + (r - q)\,S\,V_S + \kappa(\theta - v)\,V_v
        + \tfrac{1}{2} v S^2\, V_{SS} + \rho\,\xi\,v S\, V_{Sv}
        + \tfrac{1}{2} \xi^2 v\, V_{vv} = r\,V
    \end{equation*}
    \vspace{-0.9em}
  \end{block}
  \begin{itemize}
    \item Every symbol except $V$ is a known number before the solve
      starts. $r$ and $q$ come from the market, and $\kappa$, $\theta$,
      $\xi$, $\rho$ and the starting variance are the model's five
      parameters. The strike sits in the expiry payoff, not in the
      equation.
    \item The unknown is the option price $V(S, v, t)$, a whole function
      with one value per scenario, and the grid is that function written
      out cell by cell.
  \end{itemize}
  \sigcallout{The reported price is $V$ evaluated at today's share price
    and variance.}
\end{frame}

\begin{frame}{The solver}{A grid over share price and variance, solved backwards from expiry}
  \begin{sigcolumns}
    \begin{sigcolumn}
      \begin{itemize}
        \item Share price runs across the grid and variance runs down it.
          Each cell holds the option price $V$ for that pair of inputs.
        \item At expiry every cell is known exactly, so the solver starts
          there and steps \sigemph{backwards} to today.
        \item Each step rebuilds every cell from its nine neighbours,
          reading one buffer and writing another.
      \end{itemize}
    \end{sigcolumn}
    \begin{sigcolumn}
      \sigfigure{stills/still_timevalue.png}{Mid-solve. Colour is the
        option's value above the payoff.}
    \end{sigcolumn}
  \end{sigcolumns}
  \sigcallout{A full production solve is about a trillion cell updates,
    which is why the speed of this one kernel matters.}
\end{frame}

% --------------------------------------------------------------------------
% Validation. Nothing gets benchmarked until it is right.
% --------------------------------------------------------------------------

\begin{frame}{Validation}{Two properties the solver must have, run as automated tests}
  \begin{sigcolumns}
    \begin{sigcolumn}
      \begin{itemize}
        \item With $\xi = 0$ and $\rho = 0$, Heston reduces to
          Black-Scholes, which has a closed form. The solver runs the
          full stencil unchanged and must reproduce the closed form.
          Measured error 1.16e-3.
        \item Put-call parity is a model-free identity between call and
          put prices. The residual is 5.4e-6 and is unchanged while the
          price varies by thirty dollars across parameter settings.
      \end{itemize}
    \end{sigcolumn}
    \begin{sigcolumn}
      \sigfigure{stills/still_collapse.png}{The $\xi = 0$ run matching
        closed-form Black-Scholes at every step of the solve.}
    \end{sigcolumn}
  \end{sigcolumns}
\end{frame}

\begin{frame}{Three independent methods}{Grid, Fourier and Monte Carlo share no code}
  \begin{itemize}
    \item Two independent methods, sharing no code with the grid solver,
      price the same contract.
  \end{itemize}
  \sigtable{lr}{\textbf{Method} & \textbf{Price}}{%
    this solver, on the grid & \$196.1684 \\
    Fourier inversion, semi-analytic & \$196.1692 \\
    Monte Carlo, four hundred thousand paths & \$196.13 $\pm$ \$0.55 \\}
  \sigcallout{The grid and Fourier prices agree to \sigemph{four parts in a
    million}, and the Monte Carlo interval brackets both.}
\end{frame}

\begin{frame}{Convergence}{Error under grid refinement}
  \begin{center}
    \sigfigure[0.66\linewidth]{../results/convergence.png}{Error against the
      Fourier price as the grid refines.}
  \end{center}
  \sigcallout{The observed order along the refinement ladder is about two,
    and the error then lands on a discretisation floor of about \$0.014.}
\end{frame}

\begin{frame}{The test gate}{No build is benchmarked unless its tests have just passed}
  \begin{center}
    \sigfigure[0.8\linewidth]{../results/code_shots/shot_iron_rule.png}{}
  \end{center}
  \sigcallout{Every job rebuilds the binary and re-runs the full test suite
    before timing anything. A build whose tests fail exits without timing,
    so no wrong version was ever benchmarked.}
\end{frame}

% --------------------------------------------------------------------------
% Optimisation. The ladder, its one big rung, and the controls.
% --------------------------------------------------------------------------

\begin{frame}{How the optimisation was structured}{Seven solvers, one technique per level, two negative controls}
  \begin{itemize}
    \item Every level is a full copy of the level below with exactly
      \sigemph{one transformation} added, so a difference between
      neighbours measures one technique.
    \item Level 0 adds nothing, to confirm the harness itself costs
      nothing.
    \item Two \sigemph{negative controls} each remove a property the
      baseline already has, and are timed the same way.
    \item Every level must match the baseline's answer before it is
      benchmarked.
  \end{itemize}
  \sigcallout{A single before and after gives one ratio without
    attribution. The ladder attributes the speedup to individual
    techniques.}
\end{frame}

% One frame per rung. Before on the left, after on the right, and the
% measured effect in the screenshot's own caption, so the number can never
% drift from the code it belongs to.

\begin{frame}{Level 0, harness control}{Adds no technique}
  \sigrungfigure
  \begin{center}
    \sigfigure{../results/code_shots/shot_L0_anchor.png}{}
  \end{center}
  \sigrungnotes{This level runs the baseline's exact arithmetic through
    the ladder's kernel selection, changing nothing else.}{If selecting a
    kernel through a function pointer cost anything, every rung above would
    be measuring the harness rather than its technique.}
  \sigcallout{The baseline's arithmetic, line for line, behind the
    ladder's dispatch. Throughput was unchanged, so later differences
    measure the techniques themselves.}
\end{frame}

\begin{frame}{Level 1, hoisting}{Loop-invariant computation moved outside the loops}
  \sigrungfigure
  \begin{center}
    \sigfigure[0.92\linewidth]{../results/code_shots/shot_L1_hoisting.png}{}
  \end{center}
  \sigrungnotes{Values that do not change inside a loop, such as grid
    spacings, their reciprocals and per-row coefficients, are computed once
    outside it and given names.}{The intended win is fewer repeated
    operations per cell, if the compiler has not already removed them.}
  \sigcallout{No measurable change. GCC at -O2 had already hoisted all of
    this, and this level demonstrates that.}
\end{frame}

\begin{frame}{Level 2, strength reduction}{A transformation the compiler is not permitted to make}
  \sigrungfigure
  \begin{center}
    \sigfigure[0.97\linewidth]{../results/code_shots/shot_L2_strength_reduction.png}{}
  \end{center}
  \sigrungnotes{Each division by a grid constant becomes a multiplication
    by its reciprocal, computed once before the loops.}{The inner loop had
    five divisions per cell, and division is the slowest arithmetic the core
    does.}
  \sigcallout{$x/(2\,ds)$ and $x \cdot (1/(2\,ds))$ round differently in
    the last bits, so -O2 must leave the divisions alone. A divide costs
    5 to 14 cycles and does not pipeline, a multiply costs about one.}
\end{frame}

\begin{frame}{Level 3, traversal order}{The inner loop follows the memory layout}
  \sigrungfigure
  \begin{center}
    \sigfigure[0.92\linewidth]{../results/code_shots/shot_L3_order.png}{}
  \end{center}
  \sigrungnotes{The loop nest is arranged so the inner loop visits cells
    in the exact order they sit in memory.}{Sequential access lets the cache
    and prefetcher stream whole lines, while a strided walk wastes most of
    every line it fetches.}
  \sigcallout{No change, by construction. The baseline already traverses
    memory in layout order.}
\end{frame}

\begin{frame}{Level 4, loop splitting}{A separate loop for the boundary row}
  \sigrungfigure
  \begin{center}
    \sigfigure{../results/code_shots/shot_L4_split.png}{}
  \end{center}
  \sigrungnotes{The $v=0$ boundary row gets its own loop, instead of an if
    test on every cell of the main loop.}{A branch in every cell can
    mispredict and stall the pipeline, so splitting keeps the main loop free
    of branches.}
  \sigcallout{No change, by construction. The baseline already splits the
    boundary row into its own loop, and the fused version is one of the
    negative controls.}
\end{frame}

\begin{frame}{Level 5, induction variables}{Row pointers instead of index arithmetic}
  \sigrungfigure
  \begin{center}
    \sigfigure{../results/code_shots/shot_L5_pointers.png}{}
  \end{center}
  \sigrungnotes{Each row is walked with pointers that advance one step per
    cell, instead of recomputing row times width plus column for every
    neighbour.}{About ten integer additions per cell fold into the address
    arithmetic the hardware does anyway.}
  \sigcallout{A small gain, the only one after level 2. About ten integer
    additions per cell fold into the hardware's addressing modes.}
\end{frame}

\begin{frame}{Level 6, unrolling}{Four cells per loop iteration}
  \sigrungfigure
  \begin{center}
    \sigfigure{../results/code_shots/shot_L6_unroll.png}{}
  \end{center}
  \sigrungnotes{The inner loop handles four cells per iteration instead
    of one.}{That means one loop test per four cells, and more independent
    arithmetic in flight to hide instruction latencies.}
  \sigcallout{No measurable gain. The compiler already unrolls counted
    loops at -O2, so manual unrolling duplicates its work.}
\end{frame}

\begin{frame}{The ladder, measured}{Only level 2 changed throughput}
  \begin{center}
    \sigfigure[0.9\linewidth]{../results/bench_ladder_hugepage.png}{}
  \end{center}
  \sigstats{52.6/baseline Mcell per s, 131.2/after strength reduction,
    134.8/level 6, 2.56x/overall speedup}
\end{frame}

\begin{frame}{The negative controls}{Each control removes a property the baseline already had}
  \begin{center}
    \sigfigure{../results/code_shots/shot_ctl_order.png}{}
  \end{center}
  \begin{itemize}
    \item Swapping the loops makes the inner iteration jump a whole row
      per step. Same arithmetic, \sigemph{2.1x slower}. That measures the
      value of the baseline's traversal order.
    \item Fusing the boundary row back into the interior loop has no
      measurable cost. The branch predictor handles the per-cell branch.
  \end{itemize}
\end{frame}

\begin{frame}{What the speedup cost}{Fifteen units in the last place of the answer}
  \begin{center}
    \sigfigure[0.7\linewidth]{../results/code_shots/shot_optmatches.png}{}
  \end{center}
  \sigcallout{The whole speedup cost 3.3e-15 relative error, and an
    automated gate holds every level and both controls to it.}
\end{frame}

% --------------------------------------------------------------------------
% Benchmarking. The protocol first, then what it measured.
% --------------------------------------------------------------------------

\begin{frame}{The benchmark protocol}{How every number in this talk was produced}
  \sigtable{ll}{\textbf{Step} & \textbf{Rule}}{%
    where jobs run & a reserved rangpur compute node, r730-2, via sbatch \\
    before every sweep & the binary rebuilds and the full test suite must re-pass \\
    repetitions & five per point, the median is reported \\
    provenance & every CSV stamps host, job id, compiler, flags, CPU, caches \\
    the allocator & pinned, so the page regime is a choice, not an accident \\}
  \sigcallout{The plot scripts reject any CSV without its provenance
    header, so no laptop measurement can enter a figure.}
\end{frame}

\begin{frame}{Scaling with problem size}{One to sixty-four mebibytes of working set}
  \begin{center}
    \sigfigure[0.82\linewidth]{../results/bench_scaling.png}{}
  \end{center}
  \sigcallout{The first run of this sweep confounded page size with cache.
    This is the re-run with the allocator pinned, and the speedup holds at
    2.5x to 3.5x across grids.}
\end{frame}

\begin{frame}{The roofline}{Testing the memory-bound assumption}
  \begin{center}
    \sigfigure[0.75\linewidth]{../results/roofline.png}{}
  \end{center}
  \sigcallout{The solver was assumed to be memory-bound. It is not.
    Optimisation moved it from division-bound to compute-bound, with
    level 6 at \sigemph{82 percent} of the measured scalar peak.}
\end{frame}

\begin{frame}{The page-size effect}{A 1.66x change from memory pages alone}
  \vspace*{-0.4em}
  \sigtable{lrrr}{\textbf{Solver} & \textbf{4 KiB pages} & \textbf{2 MiB pages}
      & \textbf{From pages alone}}{%
    baseline & 44.8 & 52.6 & 1.17x \\
    level 2 & 80.0 & 131.2 & 1.64x \\
    \rowcolor{sigblue!12}
    level 6 & \textbf{81.1} & \textbf{134.8} & \textbf{1.66x} \\}
  \begin{itemize}
    \item Million cell updates per second. Same binary, same node, and
      not a code change at all.
    \item Sixteen mebibytes is eight huge pages or four thousand small
      ones, against a translation buffer with a thousand entries.
  \end{itemize}
  \sigcallout{It was found because the median of five repetitions was
    averaging two different memory regimes.}
\end{frame}

% --------------------------------------------------------------------------
% Reflection.
% --------------------------------------------------------------------------

\begin{frame}{Reflection}{Four questions the spec asks}
  \small
  \begin{block}{The main conclusion}
    At -O2 the compiler has already done the textbook transformations, so
    the hand optimisations that pay are the ones it is not allowed to do.
  \end{block}
  \begin{exampleblock}{The most surprising result}
    Page size changed throughput more than five of the six code techniques
    combined, without any code change.
  \end{exampleblock}
  \begin{alertblock}{What I would do differently}
    Reserve the node exclusively from the first job. An early shared-node
    measurement was real, but it answered a different question.
  \end{alertblock}
  \begin{block}{Milestone 2}
    The kernel is now compute-bound, so SIMD first, then OpenMP across
    rows. The two-buffer design already makes every cell independent
    within a step.
  \end{block}
\end{frame}

\begin{frame}[plain]
  \vfill
  \centering
  {\Large\color{sigink}\bfseries A 2.56x serial speedup, validated at
    every level}
  \par\vspace{0.6em}
  {\color{sigslate} COSC3500 Milestone 1, 48008361}
  \vfill
\end{frame}

% --------------------------------------------------------------------------
% Backup frames, outside the numbered sequence. Interview ammunition.
% --------------------------------------------------------------------------

\appendix
\sigbackupbegin

\begin{frame}{The full ladder}{Sweep E, huge pages, median of five}
  \sigtable{llrr}{\textbf{Level} & \textbf{Adds} & \textbf{Mcell per s}
      & \textbf{Step gain}}{%
    baseline & the reference solver & 52.6 & \\
    level 0 & nothing, the harness control & 52.6 & $-0.1$\,\% \\
    level 1 & hoisting, tables, CSE & 52.7 & $+0.3$\,\% \\
    \rowcolor{sigblue!12}
    level 2 & \textbf{strength reduction} & \textbf{131.2} & $\bm{+148.7}$\,\% \\
    level 3 & traversal order & 131.5 & $+0.3$\,\% \\
    level 4 & loop splitting & 131.3 & $-0.2$\,\% \\
    level 5 & induction variables & 134.2 & $+2.2$\,\% \\
    level 6 & unroll by four & 134.8 & $+0.5$\,\% \\
    \midrule
    ctl-order & loops swapped on purpose & {\color{sigred} 63.6} & 2.1x slower \\
    ctl-branch & boundary branch fused back & 135.3 & free \\}
\end{frame}

\begin{frame}{Modelled cache traffic}{Why the swapped loop order is slower}
  \begin{center}
    \sigfigure[0.86\linewidth]{stills/still_cache.png}{A modelled LRU cache
      under both traversal orders. Row-major moves 8.6 bytes per cell
      update, the swapped order moves 40.8.}
  \end{center}
\end{frame}

\sigbackupend

\end{document}
